\documentclass{article}

\textwidth=6.5in
\textheight=8.5in
\voffset=-54pt
\oddsidemargin=0in
\evensidemargin=0in
\setlength{\parskip}{8pt}
\setlength{\parindent}{0pt}

\usepackage{amsmath, amssymb}
\usepackage{ifthen}

\usepackage[inline]{enumitem}
\setlist{topsep=0pt, parsep=8pt}

\usepackage[rgb,table]{xcolor}\convertcolorsUfalse
\usepackage{graphicx}

% change hyperref options
\PassOptionsToPackage{hyphens}{url}
\usepackage[bookmarks,colorlinks,linkcolor=black,citecolor=blue,pdfstartview=FitH,urlcolor=blue]{hyperref}
\hypersetup{pdfpagemode=UseNone}
\usepackage{bookmark}

\usepackage{graphicx}
\usepackage{multicol}
\usepackage{framed}
\usepackage{array}

\usepackage{icomma}

\usepackage{soul}

\usepackage{pgfplots}
\pgfplotsset{compat=1.18}  
\pgfplotscreateplotcyclelist{mycolor}{black, blue, red, green!70!black, magenta, purple, cyan, orange }  
\pgfplotsset{
  disabledatascaling,
  cycle list=tikzcolor, 
  axis lines=middle,
  every axis plot/.style={no marks, thick},
  every axis/.style={
    enlarge x limits=false,
    enlarge y limits=auto,
    xlabel style={at={(current axis.right of origin)},anchor=west},
    ylabel style={at={(current axis.above origin)},anchor=south},
    % extra x and y ticks for asymptotes
    extra x tick style={grid=major, grid style={dashed,black}},
    extra y tick style={grid=major, grid style={dashed,black}},
    /pgf/number format/1000 sep={},
  },    
  tick label style = {font=\footnotesize, fill=\currentbgcolor, inner sep=0pt},    
  xticklabel shift=1pt,    
  scaled ticks=false,
  scaled y ticks=false,
  unbounded coords=jump,
  samples=101,
  minor tick length=0,
  scatter/classes={
    open={mark=*,draw=black,fill=white, mark size=2, thin},
    closed={mark=*,draw=black,fill=black, mark size=2, thin}
  },
  width=3in,
}
\pgfplotsset{open/.style={only marks, mark=*, fill=white, thin, mark size=2}}
\pgfplotsset{closed/.style={only marks, mark=*, draw=black, fill=black,
    thin, mark size=2}}
\pgfplotsset{labeled/.style={nodes near coords, only marks, mark=*, draw=black, fill=black, thin, mark size=2, point meta=explicit symbolic}}

\pgfplotsset{gridgraph/.style={clip=false, axis line style={shorten
      >=-8pt}, xlabel style={xshift=8pt}, ylabel style={yshift=8pt},
    enlargelimits=false, grid=major, tick style={black}}} 

\usepackage{tikz}
\usetikzlibrary{
  matrix,%
  % enables the snake paths
  decorations.pathmorphing,%
  % enables bigger arrows
  decorations.markings,%
  decorations.pathreplacing,%
  decorations.text,%
  calc,%
  patterns,%
  shapes.geometric,%
  arrows,
  decorations.shapes,
  positioning,
  plotmarks,
  intersections
}
\tikzset{>=latex} % sets default arrow type in tikz
\tikzset{font=\normalsize}
\tikzset{mark size=1.5pt, mark options=thin}
\tikzset{pin distance=4pt,
  every pin edge/.style={<-, thin, shorten <= -2pt}}

\interfootnotelinepenalty=10000
\renewcommand{\thefootnote}{(\arabic{footnote})}

\usepackage{multirow, bigdelim}

% change to \usepackage[sols]{handout} to see solutions
\usepackage[]{handout}

\newcommand\iscurrentenumi[1]{%
  \ifthenelse{\equal{#1}{\theenumi}}{}{#1}}
\setenumerate[1]{ref=\#\arabic*}
\setenumerate[2]{ref=\protect\iscurrentenumi{\theenumi}(\alph*)}
\newcommand\enumiiprefix[2]{%
  \ifthenelse{{\equal{#1}{\theenumi}}\AND{\equal{#2}{(\alph{enumii})}}}{}{}%
  \ifthenelse{{\equal{#1}{\theenumi}}\AND{\NOT{\equal{#2}{(\alph{enumii})}}}}{#2}{}%
  \ifthenelse{\equal{#1}{\theenumi}}{}{#1#2}%
}
\setenumerate[3]{ref=\protect\enumiiprefix{\theenumi}{(\alph{enumii})}\roman*}

\newlength{\tipmargin}
\makeatletter

% underline links               
\let\@href=\href
\renewcommand{\href}[2]{\@href{#1}{\ul{#2}}}

\newdimen\SOUL@dimen %new
\def\SOUL@ulunderline#1{{%
    \setbox\z@\hbox{#1}%
    % \dimen@=\wd\z@
    \SOUL@dimen=\wd\z@ %new
    \dimen@i=\SOUL@uloverlap
    \advance\SOUL@dimen2\dimen@i %\dimen@ exchanged too
    \rlap{%
      \null
      \kern-\dimen@i
      % \SOUL@ulcolor{\SOUL@ulleaders\hskip\dimen@}%
      \SOUL@ulcolor{\SOUL@ulleaders\hskip\SOUL@dimen}% new
    }%
    \unhcopy\z@
  }}

\renewenvironment{framed}{
  \setlength{\tipmargin}{\textwidth-\linewidth}
  \advance\@totalleftmargin-\tipmargin
  \def\FrameCommand{\hspace{\tipmargin}\fbox}%
  \advance\linewidth\tipmargin
  \MakeFramed {\advance\hsize-\width \FrameRestore}%
}
{\endMakeFramed}

\makeatother

\definecolor{purple}{rgb}{0.5,0,1.0}

\usepackage{fancyhdr}
\lhead{\textcolor{white}{This content is protected and may not be shared, uploaded, or distributed.}}
\rhead{}
\rfoot{\vspace*{\fill}\footnotesize\textcopyright{} \emph{Harvard Math Ma}}
\renewcommand{\headrulewidth}{0pt}
\pagestyle{fancy}
\begin{document}

\htitle{Functions and their Derivatives}

\begin{problemonly}
  \begin{framed}

You should be able to:
    \begin{itemize}[topsep=0pt,partopsep=0pt,parsep=8pt,itemsep=0pt]

    \item Explain how concavity can be characterized in terms of tangent lines.

    \item Explain what $f'$ and $f''$ tell us about the graph of $f$. (For example, if $f'$ is positive, what does that say about $f$? If $f'$ is increasing, what does that say about $f$?)

    \item Explain how displacement, velocity, and acceleration are related in terms of derivatives.

    \end{itemize}
  \end{framed}
\end{problemonly}

\begin{enumerate}
\item  \label{interpret-derivatives-cupcakes}

  \note{This is just practice interpreting the derivative, and a chance to emphasize that the estimates we make using the derivative are really just estimates.}

  \emph{Warm up.} Sara, the owner of a cupcake truck, has found that the number of cupcakes she sells per day depends on the price she charges. Let $C(p)$ be the number of cupcakes she sells when the price of a cupcake is $p$ cents. Suppose Sara has found that $C(300) = 550$ and $C'(300) = -2$. Using this information, which of the following can you find exactly? Which can you reasonably approximate?
  \begin{enumerate}
  \item
    \begin{problem}[\vfill]
      The number of cupcakes Sara will sell tomorrow.
    \end{problem}

    \begin{solution}
      We can neither find this number exactly nor reasonably approximate it. We do not have information about
      how the number of cupcakes Sara sells is related to time. 
    \end{solution}

  \item
    \begin{problem}[\vfill]
      The number of cupcakes Sara will sell in a day if she charges \$5 per cupcake.
    \end{problem}

    \begin{solution}
      We can't find this number exactly with the information given. We could attempt to approximate it
      using the information about the number of cupcakes and the instantaneous change when the price is \$3, but
      but our approximation may not be very close to the actual value, given that a price of \$5 is
      a relatively large change from a price of \$3.
    \end{solution}

  \item
    \begin{problem}[\vfill]
      The number of cupcakes Sara will sell in a day if she charges \$2.95 per cupcake.
    \end{problem}

    \begin{solution}
      We can't find this number exactly. We can approximate it using the information given, and since \$2.95
      is a relatively small change from a price of \$3, we might expect our approximation to be fairly accurate.
    \end{solution}

  \end{enumerate}

\item 
  \begin{problem}[\nstretch{1.5}]
    Previously, we described the relationship between concavity and secant lines. Now, let's think about tangent lines! 
    \begin{itemize}

    \item When a curve is concave up, its tangent lines lie \fitb{1in}{below} the curve.

    \item When a curve is concave down, its tangent lines lie \fitb{1in}{above} the curve.

    \end{itemize}
  \end{problem}

\item  \label{68-bus-route-v3}

  \note{This problem has a few goals: 
    \begin{itemize}
    \item On \href{https://people.math.harvard.edu/~jjchen/mathma/docs/homework/PS10.html}{Problem Set 10}, students did a couple examples of sketching $f'$ from a graph of $f$. This problem has them do a couple of more complex examples, to help us describe the relationships between features of $f$ and features of $f'$.

    \item Review position and velocity vs.\ speed.

    \item This problem also introduces a second derivative in context and gives students the key example of position, velocity, acceleration as an example of a function, its first derivative, and its second derivative.
    \end{itemize}
  }  

  \emph{Velocity and speed, average and instantaneous.} The \#68 bus runs
  along a straight road between Harvard Square and Kendall Square.  The
  first bus leaves Harvard Square at 6:35 am.  Let $b(t)$ be the bus's
  distance in miles from Harvard Square $t$ minutes after 6:35 am.  Here
  is the graph of $b(t)$:

  \parbox[t]{2.75in}{
    \vspace{0pt}
    \begin{tikzpicture}
      \begin{axis}[gridgraph]
        \addplot+[domain=0:22]{3.8846196535450512 - 3.8846196535450512/e^(0.43406647076875754*x) - 0.17599959491106315*x - (1.5101835483821913*x)/e^(0.43406647076875754*x) - (0.28958751199782856*x^2)/e^(0.43406647076875754*x)};      
      \end{axis}
    \end{tikzpicture}
  }
  \begin{minipage}[t]{1.0\linewidth-2.85in}
    \begin{enumerate}
    \item
      \begin{problem}[0.35in]
        How many miles long is the bus route?
      \end{problem}

      \begin{solution}
        1.5. 
      \end{solution}

    \item
      \begin{problem}[0.35in]
        At roughly what time does the bus reach Kendall Square?
      \end{problem}

      \begin{solution}
        The bus reaches Kendall Square at $t = 10$, which is 6:45 am.  
      \end{solution}

    \item
      \begin{problem}
        At roughly what time does it return to Harvard Square?
      \end{problem}

      \begin{solution}
        The bus returns to Harvard Square around $t = 22$, which is
        around 6:57 am.
      \end{solution}

    \end{enumerate}

  \end{minipage}

  \begin{enumerate}[start=4]
  \item \label{68-bus-route-velocity}

    \note{Here's an applet for the velocity: \url{https://www.geogebra.org/m/n3arwaew}}

    \begin{problem}[\vfill]
      Sketch a graph showing the bus's velocity $t$ minutes after 6:35 am
      and another showing the bus's speed $t$ minutes after 6:35 am. 
    \end{problem}

    \begin{solution}
      Since velocity is the derivative of position with respect to time, we want to sketch
      the derivative of the function given.
      \begin{center}
        \begin{tikzpicture}
          \begin{axis}[ymax=0.3, ymin=-0.2, ytick=\empty]
            \addplot+[domain=0:22]{-0.17599959491106315 + 0.17599959491106315/e^(0.43406647076875754*x) + (0.07634501906363979*x)/e^(0.43406647076875754*x) + (0.12570022931160266*x^2)/e^(0.43406647076875754*x)};
          \end{axis}
        \end{tikzpicture}
      \end{center}
      If you're not sure why the graph above makes sense, check out \href{https://www.geogebra.org/m/n3arwaew}{this applet}. Instantaneous speed is the absolute value of instantaneous velocity, so the graph looks like this.
      \begin{center}
        \begin{tikzpicture}
          \begin{axis}[ymax=0.3, ymin=-0.2, ytick=\empty]
            \addplot+[domain=0:22, samples=111]{abs(-0.17599959491106315 + 0.17599959491106315/e^(0.43406647076875754*x) + (0.07634501906363979*x)/e^(0.43406647076875754*x) + (0.12570022931160266*x^2)/e^(0.43406647076875754*x))};
          \end{axis}
        \end{tikzpicture}
      \end{center}      
    \end{solution}

  \item

    \note{This and the next part are just meant to remind students that average speed isn't necessarily the absolute value of average velocity.}

    \begin{problem}[0.1in]
      What is the bus's average velocity over this round trip?
    \end{problem}

    \begin{solution}
      The bus's average velocity is zero since the change in the bus's position is zero.
    \end{solution}

  \item
    \begin{problem}[0.1in]
      What is the bus's average speed over this round trip?
    \end{problem}

    \begin{solution}
      The bus went 3 miles in 25 minutes, so the bus's average speed is
      $\frac{3}{25} = 0.12$ miles per minute (or 7.2 miles per hour).
    \end{solution}

  \item

    \note{I'll see if students have an idea of how acceleration relates to velocity, and then I'll introduce the notation $s''$ and the term ``second derivative''.}

    \begin{problem}[0.5in]
      At roughly what times does the bus have positive acceleration? At roughly what times does it have negative acceleration?
    \end{problem}

    \begin{solution}
      By definition, acceleration is the instantaneous rate of change of velocity with respect to time. In other words, acceleration is the derivative of velocity with respect to time. 
      
      Because of this, we know that the acceleration is positive when the bus's velocity is increasing 
      and negative when the velocity is decreasing. So, the acceleration is positive on roughly $(0, 4)$ 
      and negative on roughly $(4, 22)$. We can also see that acceleration is positive when the graph of $s(t)$ 
      is concave up and negative when the graph of $s(t)$ is concave down. 

      Since velocity itself is the derivative of position
      with respect to time, acceleration is a derivative of a derivative. We call this a \textbf{second derivative}.
    \end{solution}

    \pspace{\newpage}

  \item

    \note{\url{https://www.geogebra.org/m/xub6dxwp}}

    The picture below is the actual graph of $v(t)$, the bus's velocity $t$ minutes after 6:35 am. (Does your answer to \ref{68-bus-route-velocity} match this qualitatively?) Sketch a graph of the bus's acceleration $t$ minutes after 6:35 am.

    \begin{tikzpicture}
      \begin{axis}[ytick=\empty]
        \addplot+[domain=0:22]{-0.17599959491106315 + 0.17599959491106315/e^(0.43406647076875754*x) + (0.07634501906363979*x)/e^(0.43406647076875754*x) + (0.12570022931160266*x^2)/e^(0.43406647076875754*x)};
      \end{axis}
    \end{tikzpicture}

    \begin{solution}
      Acceleration is the derivative of velocity with respect to time. Here's the graph:
      \begin{center}
        \begin{tikzpicture}
          \begin{axis}[ytick=\empty] \addplot+[domain=0:22]{-0.00005050395613637421/e^(0.43406647076875754*x) + (0.21826164563747769*x)/e^(0.43406647076875754*x) - (0.05456225491211089*x^2)/e^(0.43406647076875754*x)};
          \end{axis}
        \end{tikzpicture}
      \end{center}
      Here's \href{https://www.geogebra.org/m/xub6dxwp}{an applet that demonstrates why}. 
    \end{solution}

    \pspace{\vfill}

  \end{enumerate}

\end{enumerate}
\note{I will first just have them focus on the relationship between $f'$ and $f$. Once they're happy with that, then I'll point out that, since $f''$ is the derivative of $f'$, we have the same relationships between $f''$ and $f'$. So, this means that $f''$ tells us about concavity.}

\emph{Summarize your observations in the table below.} (If there is no
useful property, just leave the space blank.)

\begin{center}
  \renewcommand{\arraystretch}{1.75}
  \begin{tabular}{| p{1.25in} | p{1.25in} | p{1.25in} |}
    \hline
    $f$ & $f'$ & \solonly{$\bm{f''}$} \\
    \hline\hline
    positive & & \\
    \hline
    negative & & \\
    \hline\hline
    increasing $\nearrow$ & \solonly{\textbf{positive}} & \\
    \hline
    decreasing $\searrow$ & \solonly{\textbf{negative}} & \\
    \hline\hline     
    concave up \tikz[scale=0.075] \draw (-2, 4) parabola bend (0, 0) (2, 4);  & \solonly{\textbf{increasing} $\bm{\nearrow}$} & \solonly{\textbf{positive}} \\
    \hline    
    concave down \tikz[scale=0.075] \draw (-2, -4) parabola bend (0, 0) (2, -4); & \solonly{\textbf{decreasing} $\bm{\searrow}$} & \solonly{\textbf{negative}} \\
    \hline
  \end{tabular}  
  \renewcommand{\arraystretch}{1}
\end{center}

\pspace{\newpage}

\begin{enumerate}[resume]
\item  \label{learning-vocabulary}

  Seulgi is trying to memorize a set of vocabulary words for her French class. Let $W(t)$ be the number of words she knows after $t$ hours of study. She learns new words most quickly at first; as time goes on, her rate of learning keeps getting slower.

  \begin{enumerate}
  \item
    \begin{problem}[0.25in]
      What does the last sentence tell you about $W'(t)$? 
    \end{problem}

    \begin{solution}
      $W'$ represents the rate of learning, so $W'$ is decreasing.
    \end{solution}

  \item
    \begin{problem}[0.25in]
      What does the last sentence tell you about $W''(t)$?      
    \end{problem}

    \begin{solution}
      If $W'$ is decreasing, we can deduce that $W''$ is negative.
    \end{solution}

  \item
    \begin{problem}[0.25in]
      What does the last sentence tell you about $W(t)$? 
    \end{problem}

    \begin{solution}
      $W$ is concave down.
    \end{solution}

  \end{enumerate}

\item  \label{interpret-trip-1}

  North Dakota's Highway 46 is supposedly the straightest road in the
  United States.  From the city of Streeter, it runs 123 miles east to
  Hickson.  Litchville is a tiny city on this highway.

  A group of friends decides to drive on the highway one day.  Let $s(t)$
  be their position at time $t$, where $t$ is measured in minutes after
  noon and the position is given in miles east of Litchville.  Here is the
  graph of $s(t)$:

  \begin{tikzpicture}
    \begin{axis}[declare function={
        f(\x) = and(0. <= \x, \x <=10) * (-1.5 + 0.07888918067226876*\x + 0.014013452290002856*\x^2 - 0.0006902370357229733*\x^3)
        + and(10 < \x, \x <=15) * (1.0291392835516033 - 0.5973651562133899*\x + 0.07339014116058672*\x^2 - 0.002394501837476374*\x^3)
        + and(15 < \x, \x <=20) * (3.157905419052831 - 0.664484013390334*\x + 0.05395577364416258*\x^2 - 0.0014313168261138993*\x^3)
        + and(20 < \x, \x <=30) * (12.289986612892015 - 1.1004984381707412*\x + 0.029066607168409475*\x^2 - 0.00023833258960512233*\x^3)
        + and(30 < \x, \x <=40) * (39.4999999999985 - 3.5999999999998895*\x + 0.10499999999999728*\x^2 - 0.0009999999999999786*\x^3)
        + and(40 < \x, \x <=45.036) * (-399.0887813408259 + 28.311260641379604*\x - 0.6682090670549302*\x^2 + 0.005238638483961406*\x^3)
        + and(45.036 < \x, \x <=60) * (67.34196489531192 - 4.03793774180702*\x + 0.07847962064747413*\x^2 - 0.0004981126607305439*\x^3)
        ;}, gridgraph, xtick={5, 10, ..., 60}, ytick={-9, -6, ..., 3}, xlabel=time
      (minutes), ylabel=position (miles), width=4in, height=2.5in]
      \addplot+[domain=0:60]{6*f(x)};
    \end{axis}
  \end{tikzpicture}

  \begin{enumerate}

  \item
    \begin{problem}[\vfill]
      When are the friends traveling east?  
    \end{problem}

    \begin{solution}
      The friends are traveling east when $t$ is in the intervals $(0, 15) \cup (30, 40) \cup (45, 60)$, so
      from 12:00 to 12:15, 12:30 to 12:40, and 12:45 to 1:00.
    \end{solution}

    \begin{problem}[\vfill]
      What did you look at to decide? 
    \end{problem}

    \begin{solution}
      Since $s(t)$ is the number of miles east of Litchville, the friends are traveling east when $s(t)$ is increasing.
    \end{solution}

  \item
    \begin{problem}[\vfill\newpage]
      Are the friends going faster at 12:05 or at 12:35?  How do you know?
    \end{problem}

    \begin{solution}
      Looking at the graph, it seems that the friends are going faster at 12:05 than they are at 12:35. This
      is because the slope of the tangent line at $t=5$ appears to be steeper than the slope of the tangent 
      line at $t=35$.
    \end{solution}

  \end{enumerate}

  Now, suppose that a different group of friends is driving on the same 
  highway, and that this is the graph of their \emph{velocity} at time
  $t$:

  \begin{tikzpicture}
    \begin{axis}[declare function={
        f(\x) = and(0. <= \x, \x <=10) * (-1.5 + 0.07888918067226876*\x + 0.014013452290002856*\x^2 - 0.0006902370357229733*\x^3)
        + and(10 < \x, \x <=15) * (1.0291392835516033 - 0.5973651562133899*\x + 0.07339014116058672*\x^2 - 0.002394501837476374*\x^3)
        + and(15 < \x, \x <=20) * (3.157905419052831 - 0.664484013390334*\x + 0.05395577364416258*\x^2 - 0.0014313168261138993*\x^3)
        + and(20 < \x, \x <=30) * (12.289986612892015 - 1.1004984381707412*\x + 0.029066607168409475*\x^2 - 0.00023833258960512233*\x^3)
        + and(30 < \x, \x <=40) * (39.4999999999985 - 3.5999999999998895*\x + 0.10499999999999728*\x^2 - 0.0009999999999999786*\x^3)
        + and(40 < \x, \x <=45.036) * (-399.0887813408259 + 28.311260641379604*\x - 0.6682090670549302*\x^2 + 0.005238638483961406*\x^3)
        + and(45.036 < \x, \x <=60) * (67.34196489531192 - 4.03793774180702*\x + 0.07847962064747413*\x^2 - 0.0004981126607305439*\x^3)
        ;}, gridgraph, xtick={5, 10, ..., 60}, ytick={-45, -30, ..., 15}, xlabel=time
      (minutes), ylabel=velocity (mph), width=4in, height=2.5in]
      \addplot+[domain=0:60]{30*f(x)};
    \end{axis}
  \end{tikzpicture}

  \begin{enumerate}[resume]

  \item
    \begin{problem}[0.25in]
      When are the friends traveling east?  
    \end{problem}

    \begin{solution}
      The friends are traveling east when $t$ is in the interval $(10, 20)$, so between 12:10 and 12:20.
    \end{solution}

    \begin{problem}[0.25in]
      What did you look at to decide?
    \end{problem}

    \begin{solution}
      Traveling east corresponds to the velocity being positive.
    \end{solution}

  \item 
    \begin{problem}[0.25in]
      Are the friends going faster at 12:30 or 12:35?  How do you know?
    \end{problem}

    \begin{solution}
      The friends are traveling faster at 12:30 than at 12:35. This is because instantaneous speed is the absolute value
      of instantaneous velocity, and based on the graph, the instantaneous speed is greater when $t=30$ than when
      $t=35$. 
    \end{solution}

  \end{enumerate}

  \pspace{\newpage}

\item 
  Let $f(x) = (x - 2)^3 + 1$.

  \begin{enumerate}
  \item
    \begin{problem}[\vfill]
      Sketch a graph of $f$.
    \end{problem}

    \begin{solution}
      \begin{center}
        \begin{tikzpicture}
          \begin{axis}[ytick=\empty]
            \addplot[domain=0:4]{(x-2)^3 +1};
          \end{axis}
        \end{tikzpicture}
      \end{center}
    \end{solution}

  \item
    \begin{problem}[\vfill]
      Based on your graph of $f$, sketch a graph of $f'$. 
    \end{problem}

    \begin{solution}
    Since $f$ is always increasing, $f'$ should be positive except at $x=2$, where there is
    a horizontal tangent. At $x=2$, $f'$ should be zero. Since $f$ is concave down when $x<2$ and
    concave up when $x>2$, $f'$ should be decreasing when $x<2$ and increasing when $x>2$. 
      \begin{center}
        \begin{tikzpicture}
          \begin{axis}[ytick=\empty]
            \addplot[domain=0:4]{3*(x-2)^2};
          \end{axis}
        \end{tikzpicture}
      \end{center}
    \end{solution}

  \item
    \begin{problem}[\stretch{1.75}]
      Use the definition of the derivative to calculate $f'(2)$.
    \end{problem}

    \begin{solution}
      \begin{align*}
        f'(2) &= \lim_{h \to 0} \frac{f(2+h) - f(2)}{h} \\
        &= \lim_{h \to 0} \frac{\bigl((2+h-2)^3 +1\bigr) - \bigl((2-2)^3 +1 \bigr)}{h} \\
        &= \lim_{h \to 0} \frac{h^3}{h} \\
        &= \lim_{h \to 0} h^2 \\
        &= \boxed{0}
      \end{align*}
    \end{solution}

  \item
    \begin{problem}[\stretch{0.25}]
      How is the value of $f'(2)$ reflected in your graph of $f$? How is it reflected in your graph of $f'$? 
    \end{problem}

    \begin{solution}
      On the graph of $f$, we can see that there is a horizontal tangent at $x=2$. On the graph of $f'$, there
      is a $x$-intercept at $x=2$. 
    \end{solution}

  \end{enumerate}

\end{enumerate}
\end{document}
